\chapter*{Abstract}
\phantomsection
\addcontentsline{toc}{chapter}{Abstract}

\begin{otherlanguage}{english}

An algorithm for fusing infrared and visible images based on mathematical morphology is proposed: a
single-scale Top-Hat with a disk and four oriented line segments. It adds the bright detail to the average
of the two images and subtracts the dark detail, both multiplied by a weight. Through particle swarm
optimization, the weight settled at the lowest allowed value, 0.30, and the radius of 25 pixels, which
gives more enhancement, was chosen. The hypothesis is that the algorithm outperforms the classical Top-Hat
in quality and in detecting objects seen by only one camera, and that it is comparable in quality to five
multiscale methods. With six no-reference metrics on twenty nighttime pairs, it outperformed the classical
Top-Hat on four, lost on one and tied on another. It achieved the best average rank of the seven methods,
though narrowly and with an order that changes with the metrics: against the multiscale methods, it gained
contrast and sharpness and lost similarity to the source images. On another dataset, with daytime and
nighttime scenes, the balanced hit rate of the detector trained on its fused images was 61.4\%. It averages the
hit rates on objects seen only by the visible camera and on those seen only by the infrared one. It ranked
fourth among eight fused inputs and was 3.0 points above the classical Top-Hat, a minimal advantage that
changed sign with the seed. The hypothesis largely holds for quality; for detection, it remains to be
confirmed with new training runs and data.

\textbf{Keywords:} Infrared and visible image fusion; Mathematical morphology; Top-Hat transform;
Particle swarm optimization; Object detection.

\end{otherlanguage}
